\scanpage{39}\mspage{32}
% archive sheet number: 38
coïncident, et deviennent \inserted{$\unsure{=|\simeq}$} particulièrement simples. (Car si $f,g:X\unsure{=|\rightrightarrows}Y$, avec $X$ cofibrant et $Y$ fibrant, alors $f\sim g$ peut se tester sur \uline{n'importe} \uline{quel} ``objet cylindre'' pour $X$, ou sur \uline{n'importe quel} ``objet chemins'' pour $Y$.)

\inserted{Dans} le cas où la structure de Quillen faible est exacte à gauche, \deleted{et} i.e. où la relation d'équi\unsure{v.} pertinente dans $\operatorname{Fl}M_{cf}$ est la Q-homotopie à droite, alors l'axiome d'exactitude est vérifié pour les carrés $W$-cartésiens, mais je doute qu'il en soit de $=$ pour les carrés $W$-cocartésiens.$^{*}$
% slanted
\margin{Il semble bien \unsure{que c'est} le cas, d'après \unsure{Baues|Barnes}, \unsure{cf annotation|cf annotations} \unsure{marginale correspondante|marginales correspondantes} à page 30}
(Si on ne fait \deleted{\unsure{l'}} une hypothèse supplémentaire, p. ex. l'exactitude à dr. de la structure de Quillen) Si on n'a que l'exactitude (mais ni à g. ni à dr.), il semble qu'on ne peut en conclure l'exactitude ni des carrés $W$-cartésiens, ni des carrés $W$-cocartésiens, \deleted{On peut} \inserted{\unsure{donc}} pas celle des suites de \unsure{cosuspension|cosuspensions} (loop \unsure{sequence|sequences}), ni des suites de suspension. On pourrait donc penser que la condition d'exactitude, à elle seule, ne sert pas à grand-
